\documentclass[10pt,aspectratio=169]{beamer}

\usepackage[utf8]{inputenc}
\usepackage{xeCJK}
\usepackage{graphicx}
\usepackage{mathtools}
\usepackage{utopia}
\usepackage{booktabs}
\usepackage{xcolor}
\usepackage{natbib}
\usepackage{hyperref}

\usetheme{CambridgeUS}
\usecolortheme{dolphin}

\definecolor{myNewColorA}{RGB}{118,193,188}
\definecolor{myNewColorB}{RGB}{106,172,150}
\definecolor{myNewColorC}{RGB}{94,150,218}
\setbeamercolor*{palette primary}{bg=myNewColorC}
\setbeamercolor*{palette secondary}{bg=myNewColorB, fg=white}
\setbeamercolor*{palette tertiary}{bg=myNewColorA, fg=white}
\setbeamercolor*{titlelike}{fg=myNewColorA}
\setbeamercolor*{title}{bg=myNewColorA, fg=white}
\setbeamercolor*{item}{fg=myNewColorA}
\setbeamercolor*{caption name}{fg=myNewColorA}
\usefonttheme{professionalfonts}
\setbeamertemplate{navigation symbols}{}
\graphicspath{{../../charts/}{../shared/figures/}}

\title[The University of Hong Kong]{Data Preparation and Descriptive Analysis}
\subtitle{JFK Airline Delay and Cancellation Operational Risk}
\author[DASE7337 Group Project]{DASE7337 Group Project}
\institute[HKU]{The University of Hong Kong}
\date[April 2026]{DASE7337 Group Presentation \\ April 2026}
\titlegraphic{
\includegraphics[height=1.5cm]{hkulogo.png}}

\setbeamerfont{title}{size=\large}
\setbeamerfont{subtitle}{size=\small}
\setbeamerfont{author}{size=\small}
\setbeamerfont{date}{size=\small}
\setbeamerfont{institute}{size=\small}

\begin{document}

\begin{frame}
  \titlepage
\end{frame}

\begin{frame}{Where This Chapter Fits}
  % Speaking cue: 30--35 seconds.
  \begin{itemize}
    \item The previous chapter covers sector context, core risk factors, SWOT, and PEST.
    \item This chapter converts BTS raw records into modeling-ready risk assets.
    \item The key question is not only which charts we produced, but which data definitions can support the next models.
  \end{itemize}
  \vspace{0.4em}
  \begin{block}{Modeling handoff}
  The airline-month records are organized into frequency variables, severity variables, and delay-cause variables for distribution fitting, state classification, and aggregate risk simulation.
  \end{block}
\end{frame}

\begin{frame}{Data Scope and Cleaning Result}
  % Speaking cue: 35--40 seconds.
  \begin{columns}[T,totalwidth=\linewidth]
    \begin{column}{0.48\linewidth}
      \begin{table}
      \small
      \begin{tabular}{lr}
      \toprule
      Metric & Value \\
      \midrule
      Raw records & 570 \\
      Clean records & 569 \\
      Airport months & 72 \\
      Airlines & 12 \\
      Arrival flights & 657,671 \\
      Delay minutes & 11,125,838 \\
      \bottomrule
      \end{tabular}
      \end{table}
    \end{column}
    \begin{column}{0.48\linewidth}
      \begin{itemize}
        \item One non-operational record was removed.
        \item Logical check: delayed arrivals must not exceed total arrivals.
        \item Airline names were standardized and BTS fields were rewritten into readable analytical names.
      \end{itemize}
    \end{column}
  \end{columns}
\end{frame}

\begin{frame}{Variable Design for Frequency--Severity Modeling}
  % Speaking cue: 40 seconds.
  \begin{columns}[T,totalwidth=\linewidth]
    \begin{column}{0.33\linewidth}
      \textbf{Frequency}
      \begin{itemize}
        \item Arrival flights
        \item 15+ minute delays
        \item Cancellations and diversions
      \end{itemize}
    \end{column}
    \begin{column}{0.33\linewidth}
      \textbf{Severity}
      \begin{itemize}
        \item Total delay minutes
        \item Average minutes per delay
        \item Cause-specific minutes
      \end{itemize}
    \end{column}
    \begin{column}{0.33\linewidth}
      \textbf{Risk Blocks}
      \begin{itemize}
        \item Airline and Late Aircraft
        \item NAS
        \item Weather and Security
      \end{itemize}
    \end{column}
  \end{columns}
  \vspace{0.8em}
  \centering
  \large
  The cleaned structure lets later models ask separately: how often, and how severe?
\end{frame}

\begin{frame}{Monthly Trend: Delays and Cancellations Differ}
  % Speaking cue: 45 seconds.
  \centering
  \includegraphics[width=0.88\linewidth]{chart_1_monthly_trend.png}
  \vspace{-0.3em}
  \begin{itemize}
    \item Peak delay month: July 2023, with 3,657 delayed arrivals and a 32.69\% delay rate.
    \item Peak cancellation month: March 2020, with 1,900 cancellations during the early pandemic shock.
  \end{itemize}
\end{frame}

\begin{frame}{Delay-Cause Structure: Internal Chain Dominates}
  % Speaking cue: 45 seconds.
  \centering
  \includegraphics[width=0.86\linewidth]{chart_2_delay_cause_breakdown.png}
  \vfill
  \begin{minipage}{0.86\linewidth}
    \centering\small
    Airline and Late Aircraft jointly contribute about 71.3\% of delay minutes, so daily operational risk is driven mainly by the internal operating chain and system congestion, not only by weather.
  \end{minipage}
\end{frame}

\begin{frame}{Airline Risk Profile and Dashboard Handoff}
  % Speaking cue: 45--50 seconds.
  \begin{columns}[T,totalwidth=\linewidth]
    \begin{column}{0.60\linewidth}
      \includegraphics[width=\linewidth]{chart_3_airline_risk_profile.png}
    \end{column}
    \begin{column}{0.36\linewidth}
      \begin{itemize}
        \item JetBlue has the largest total delay minutes.
        \item Hawaiian has the highest delay rate.
        \item Envoy Air has the highest average delay severity.
        \item The offline dashboard supports filtering, interactive charts, and a readable data table.
      \end{itemize}
      \vspace{0.3em}
      \href{../../dashboard/jfk_operational_risk_dashboard.html}{Open project dashboard}
    \end{column}
  \end{columns}
\end{frame}

\begin{frame}{Why This Matters for the Next Models}
  % Speaking cue: 35--40 seconds.
  \begin{enumerate}
    \item The cleaned core dataset gives clear inputs for frequency and severity models.
    \item Cause categories translate BTS fields into operational risk blocks.
    \item The monthly time series supports normal/disrupted state classification and Markov transition estimation.
    \item The charts and dashboard make the modeling story explainable in business language.
  \end{enumerate}
\end{frame}

\end{document}
